% ---- pg-593 （书页 581） ----

{\small
\[
\begin{array}{r@{\;}c}
 & \begin{array}{ccccc}\text{JPMorgan}&\text{Citibank}&\text{Wells Fargo}&\text{Royal Dutch Shell}&\text{ExxonMobil}\end{array}\\[4pt]
\begin{array}{r}\text{JPMorgan}\\\text{Citibank}\\\text{Wells Fargo}\\\text{Royal Dutch Shell}\\\text{ExxonMobil}\end{array} &
\left[\begin{array}{ccccc}
1&&&&\\
0.63&1&&&\\
0.51&0.57&1&&\\
0.12&0.32&0.18&1&\\
0.16&0.21&0.15&0.68&1
\end{array}\right]
\end{array}
\]
}

将样本相关系数用作相似性度量, 利用练习 12.5 中的三种方法将五种股票聚类. 画出树图并比较所得结果.

\noindent 12.8\quad 利用例 12.3 中的距离数据, 用平均连接分层方法对项目进行聚类. 画出树图. 将结果与例 12.3 和 12.5 中的结果进行比较.

\noindent 12.9\quad 一本书中词汇的“丰富性”可以通过统计用过一次的词、用过两次的词等来描述. 根据这些统计数字, 一个语言学家给出了《旧约全书》耶利米哀歌的各章之间的距离如下:

\[
\begin{array}{c@{\quad}r@{\;}c}
 & \multicolumn{2}{c}{\text{耶利米哀歌, 章}}\\[4pt]
 & \begin{array}{ccccc}1&2&3&4&5\end{array}\\[4pt]
\text{耶利米哀歌, 章} & \begin{array}{r}1\\2\\3\\4\\5\end{array} &
\left[\begin{array}{ccccc}
0&&&&\\
0.76&0&&&\\
2.97&0.80&0&&\\
4.88&4.17&0.21&0&\\
3.86&1.92&1.51&0.51&0
\end{array}\right]
\end{array}
\]

我们用讨论过的三种连接分层聚类方法对耶利米哀歌各章进行聚类. 画树图并比较结果.

\noindent 12.10\quad 用沃德方法对四个项目聚类. 这些项目在变量 $X$ 上的测量值列于下表.

\begin{mvtable}
\setlength{\tabcolsep}{8pt}%
\begin{tabular}{c|c}
\toprule
\textbf{项 目} & \textbf{测 量 值} \\
\midrule
1 & 2 \\
2 & 1 \\
3 & 5 \\
4 & 8 \\
\bottomrule
\end{tabular}
\end{mvtable}

\noindent (a) 起初, 每个项目为一个聚类, 因而我们有类 $\{1\},\{2\},\{3\},\{4\}$. 证明 ESS$=0$.

\noindent (b) 若我们联合聚类 $\{1\}$ 和 $\{2\}$, 则新聚类 $\{12\}$ 有
\[
\mathrm{ESS}_1=\sum(x_j-\bar{x})^2=(2-1.5)^2+(1-1.5)^2=0.5
\]
于是与分组 $\{12\},\{3\},\{4\}$ 相联系的 ESS 为 ESS$=0.5+0+0=0.5$. ESS 的增量(信息损失)为 $0.5-0=0.5$. 对步骤 2 中的所有可能性确定 ESS 的增量, 并将结果填入下表.

\begin{mvtable}
\setlength{\tabcolsep}{10pt}%
\begin{tabular}{ccc|c}
\toprule
\multicolumn{3}{c|}{\textbf{聚 类}} & \textbf{ESS 的增量} \\
\midrule
$\{12\}$ & $\{3\}$ & $\{4\}$ & 0.5 \\
$\{13\}$ & $\{2\}$ & $\{4\}$ & \\
$\{14\}$ & $\{2\}$ & $\{3\}$ & \\
$\{1\}$ & $\{23\}$ & $\{4\}$ & \\
$\{1\}$ & $\{24\}$ & $\{3\}$ & \\
$\{1\}$ & $\{2\}$ & $\{34\}$ & \\
\bottomrule
\end{tabular}
\end{mvtable}
